% Part: history
% Chapter: biographies
% Section: alan-turing

\documentclass[../../../include/open-logic-section]{subfiles}

\begin{document}

\olfileid{his}{bio}{tur}

\olsection{آلن تورینگ}

آلن تورینگ در ۲۳ ژوئن ۱۹۱۲ در مِیدا وِیلِ لندن زاده شد. او را پدر علوم
نظری رایانه می‌دانند. علاقهٔ تورینگ به علوم طبیعی و ریاضیات از کودکی آغاز
شد. بااین‌حال، مدرسه‌هایش در دوران کودکی بازتاب مناسبی از علایق او نداشتند
و بر ادبیات و آثار کلاسیک تأکید می‌کردند. ازاین‌رو در مدرسه عملکرد خوبی
نداشت و بسیاری از آموزگارانش او را سرزنش می‌کردند.

\olphoto{turing-alan}{Alan Turing}

تورینگ دورهٔ کارشناسی را در کینگز کالجِ کمبریج گذراند و در آنجا ریاضیات
خواند. او در سال ۱۹۳۶، برای تعریف دقیق مفهوم تابع محاسبه‌پذیر و اثبات
تصمیم‌ناپذیری مسئلهٔ تصمیم، چیزی را ساخت که امروزه «ماشین تورینگ» نامیده
می‌شود. آلونزو چرچ با اثبات همان نتیجه به‌وسیلهٔ حساب لامبدای خود زودتر به
آن دست یافته بود. بااین‌همه، مقالهٔ تورینگ با ارجاع به نتیجهٔ چرچ منتشر
شد. چرچ تورینگ را به پرینستون دعوت کرد؛ تورینگ از ۱۹۳۶ تا ۱۹۳۸ در آنجا
بود و زیر نظر چرچ دکتری گرفت.

با وجود علاقهٔ تورینگ به منطق، علایق پیشین او به علوم طبیعی همچنان برجسته
ماند. مهارت‌های عملی او در دوران خدمتش در واحد تحلیل رمز بریتانیا در
بلچلی پارک، طی جنگ جهانی دوم، به کار آمد. تورینگ از چهره‌های محوری شکستن
رمز ارتباطات نیروی دریایی آلمان، یعنی رمز انیگما، بود. تخصص او در آمار و
رمزنگاری، همراه با ورود ماشین‌های الکترونیکی، به گروه امکان داد با ساختن
دستگاه رمزگشایی‌ای به نام «بامب» رمز را بشکند. اندیشه‌هایش در ساخت
کولوسوس، نخستین رایانهٔ الکترونیکی برنامه‌پذیر جهان، نیز مؤثر بود؛ این
رایانه در بلچلی پارک برای شکستن رمز لورنتس آلمان به کار رفت.

تورینگ همجنس‌گرا بود. بااین‌حال، در سال ۱۹۴۲ از جوآن کلارک، یکی از
همکارانش در بلچلی پارک، خواستگاری کرد؛ اما بعداً نامزدی را برهم زد و به او
گفت که همجنس‌گراست. او در طول زندگی چند شریک عاطفی داشت، درحالی‌که روابط
جنسی میان مردان در آن زمان در بریتانیا جرم بود. در سال ۱۹۵۲، یکی از
دوستان شریک عاطفی وقت تورینگ به خانه‌اش دستبرد زد. تورینگ هنگام گزارش
واقعه به پلیس به داشتن رابطه‌ای همجنس‌خواهانه اقرار کرد، زیرا گمان می‌کرد
دولت در آستانهٔ قانونی‌کردن این روابط است. چنین نبود و او به «بی‌عفتی
فاحش» متهم شد. تورینگ به‌جای زندان، درمان هورمونی کاهندهٔ میل جنسی را
برگزید. پیکر بی‌جان او در ۸ ژوئن ۱۹۵۴ در پی مصرف بیش‌ازحد سیانور پیدا
شد---به احتمال زیاد خودکشی کرده بود. ملکه الیزابت دوم در سال ۲۰۱۳ او را
رسماً عفو سلطنتی کرد.

\begin{reading}
برای زندگی‌نامه‌ای جامع از آلن تورینگ، بنگرید به \citet{Hodges2014}.
زندگی و کار تورینگ الهام‌بخش نمایشنامهٔ \emph{شکستن رمز} بود که در سال
۱۹۹۶ با بازی درک جاکوبی در نقش تورینگ برای تلویزیون ساخته شد.
\emph{بازی تقلید}، فیلم نامزد جایزهٔ اسکار با بازی بندیکت کامبربچ و کیرا
نایتلی، نیز اقتباسی آزاد از زندگی آلن تورینگ و دوران حضورش در بلچلی پارک
است \citep{Imitation2014}.

\citet{Radiolab2012} چند پادکست دربارهٔ زندگی و کار تورینگ دارد. مستند
برنامهٔ هورایزن بی‌بی‌سی با عنوان \emph{زندگی و مرگ شگفت‌انگیز دکتر
تورینگ} برای تماشای برخط در دسترس است \citep{Sykes1992}.
\citep{Theelen2012} ویدیوی کوتاهی از یک ماشین تورینگ لگوییِ فعال است که
برای بزرگداشت صدمین سال تولد تورینگ در سال ۲۰۱۲ ساخته شد.

مقالهٔ اصلی تورینگ دربارهٔ ماشین‌های تورینگ و مسئلهٔ تصمیم در
\citet{Turing1937} آمده است.
\end{reading}

\end{document}
